\par\Needspace{6\baselineskip}\medskip
\hypertarget{native-102-article077}{}\pdfbookmark[2]{77.}{article-077}
\providecommand{\DailyMotionSign}{\mathord{\text{\itshape 7}}}
\begin{center}77.\end{center}
Combining what was developed in the preceding article with what we presented in Arts. 15, 16, 20, 27 and 28 for the several kinds of conic section gives us everything needed to calculate the differential changes in the geocentric position resulting from changes in the individual elements. To illustrate these rules more fully, we shall return to the example discussed above in Arts. 13, 14, 51, 63 and 65. First, following the preceding article, we shall express $\mathrm{d}l$ and $\mathrm{d}b$ in terms of $\mathrm{d}r$, $\mathrm{d}u$, $\mathrm{d}i$ and $\mathrm{d}\AscendingNode$. The calculation proceeds as follows.
\begin{center}
\begin{tabular}{@{}ll@{}}
\makebox[17em][l]{$\log\operatorname{tang}\omega$\hfill .\,.\,.\,.} & $8{,}40113$\\[.4em]
\makebox[17em][l]{$\log\cos i$\hfill .\,.\,.\,.\,.} & $9{,}98853$\\[.4em]
\hline\noalign{\vskip .3em}
\makebox[17em][l]{$\log\operatorname{tang}M$\hfill .\,.\,.} & $8{,}41260$\\[.4em]
$M=$ & $1^0\,28'\,52''$\\[.4em]
$M-u=$ & $165\quad17\quad8$\\[.4em]
\end{tabular}
\end{center}
\begin{center}
\begin{tabular}{@{}ll@{}}
\makebox[17em][l]{$\log\sin\omega$\hfill .\,.\,.\,.\,.} & $8{,}40099_n$\\[.4em]
\makebox[17em][l]{$\log\operatorname{tang}i$\hfill .\,.\,.\,.} & $9{,}36723$\\[.4em]
\hline\noalign{\vskip .3em}
\makebox[17em][l]{$\log\operatorname{tang}N$\hfill .\,.\,.} & $7{,}76822_n$\\[.4em]
$N=$ & $179^0\,39'\,50''$\\[.4em]
$N-b=$ & $186\quad1\quad45$\\[.4em]
\end{tabular}
\end{center}
\begin{center}
\begin{tabular}{@{}ll@{}}
\makebox[17em][l]{$\log\operatorname{tang}(M-u)$\hfill .\,.\,.} & $9{,}41932_n$\\[.4em]
\makebox[17em][l]{$\log\cos\omega\sin i$\hfill .\,.\,.\,.\,.} & $9{,}35562_n$\\[.4em]
\hline\noalign{\vskip .3em}
\makebox[17em][l]{$\log\operatorname{tang}P$\hfill .\,.\,.\,.\,.\,.\,.\,.} & $0{,}06370$\\[.4em]
$P=$ & $49^0\,11'\,13''$\\[.4em]
$N-b-P=$ & $136\quad50\quad32$\\[.4em]
\end{tabular}
\end{center}
\SourceLeaf{103}
\begin{center}
\begin{tabular}{@{}ll@{}}
\multicolumn{2}{c}{I$^{*}$}\\[.4em]
\makebox[17em][l]{$\log\sin(L-l)$\hfill .\,.\,.} & $9{,}72125$\\[.4em]
\makebox[17em][l]{$\log R$\hfill .\,.\,.\,.\,.\,.\,.\,.\,.\,.} & $9{,}99810$\\[.4em]
\makebox[17em][l]{$\mathrm{C.}\log\Delta'$\hfill .\,.\,.\,.\,.\,.\,.} & $9{,}92027$\\[.4em]
\hline\noalign{\vskip .3em}
\makebox[17em][l]{$({}^{*})$\hfill .\,.\,.\,.\,.\,.\,.\,.\,.\,.\,.\,.} & $9{,}63962$\\[.4em]
\makebox[17em][l]{$\mathrm{C.}\log r$\hfill .\,.\,.\,.\,.\,.\,.\,.} & $9{,}67401$\\[.4em]
\hline\noalign{\vskip .3em}
\makebox[17em][l]{$\log\left(\frac{\mathrm{d}l}{\mathrm{d}r}\right)$\hfill .\,.\,.\,.\,.\,.\,.} & $9{,}31363$\\[.4em]
\end{tabular}
\end{center}
\begin{center}
\begin{tabular}{@{}ll@{}}
\multicolumn{2}{c}{II$^{**}$}\\[.4em]
\makebox[17em][l]{$({}^{*})$\hfill .\,.\,.\,.\,.\,.\,.\,.\,.\,.\,.} & $9{,}63962$\\[.4em]
\makebox[17em][l]{$\log\operatorname{cot}(M-u)$\hfill .} & $0{,}58068_n$\\[.4em]
\hline\noalign{\vskip .3em}
\makebox[17em][l]{$\log\left(\frac{\mathrm{d}l}{\mathrm{d}u}\right)$\hfill .\,.\,.\,.\,.\,.} & $0{,}22030$\\[.4em]
\end{tabular}
\end{center}
\begin{center}
\begin{tabular}{@{}ll@{}}
\multicolumn{2}{c}{III$^{*}$}\\[.4em]
\makebox[17em][l]{$\log\cos\omega$\hfill .\,.\,.\,.\,.\,.\,.\,.\,.\,.\,.} & $9{,}99986_n$\\[.4em]
\makebox[17em][l]{$\log\operatorname{tang}b$\hfill .\,.\,.\,.\,.\,.\,.\,.\,.\,.} & $9{,}04749_n$\\[.4em]
\hline\noalign{\vskip .3em}
\makebox[17em][l]{$\log\left(\frac{\mathrm{d}l}{\mathrm{d}i}\right)$\hfill .\,.\,.\,.\,.\,.\,.\,.\,.\,.\,.} & $9{,}04735_n$\\[.4em]
\end{tabular}
\end{center}
\begin{center}
\begin{tabular}{@{}ll@{}}
\multicolumn{2}{c}{IV}\\[.4em]
\makebox[17em][l]{$\log\frac{R}{\Delta'}$\hfill .\,.\,.\,.\,.\,.\,.\,.\,.\,.} & $9{,}91837$\\[.4em]
\makebox[17em][l]{$\log\cos(L-l)$\hfill .\,.} & $9{,}92956$\\[.4em]
\hline\noalign{\vskip .3em}
\makebox[17em][l]{$({}^{**})$\hfill .\,.\,.\,.\,.\,.\,.\,.\,.\,.\,.} & $9{,}84793$\\[.4em]
\multicolumn{2}{@{}l@{}}{$=\log\left(\frac{\mathrm{d}l}{\mathrm{d}\AscendingNode}-1\right)$}\\[.4em]
\end{tabular}
\end{center}
\begin{center}
\begin{tabular}{@{}ll@{}}
\multicolumn{2}{c}{V$^{*}$}\\[.4em]
\makebox[17em][l]{$({}^{**})$\hfill .\,.\,.\,.\,.\,.\,.\,.\,.\,.} & $9{,}84793$\\[.4em]
\makebox[17em][l]{$\log\sin b\cos b$\hfill .\,.} & $9{,}04212_n$\\[.4em]
\makebox[17em][l]{$\mathrm{C.}\log r$\hfill .\,.\,.\,.\,.\,.\,.} & $9{,}67401$\\[.4em]
\hline\noalign{\vskip .3em}
\makebox[17em][l]{$\log\left(\frac{\mathrm{d}b}{\mathrm{d}r}\right)$\hfill .\,.\,.\,.\,.\,.} & $8{,}56406$\\[.4em]
\end{tabular}
\end{center}
\begin{center}
\begin{tabular}{@{}ll@{}}
\multicolumn{2}{c}{VI$^{**}$}\\[.4em]
\makebox[17em][l]{$\log\frac{r}{\Delta}$\hfill .\,.\,.\,.\,.\,.\,.\,.\,.\,.\,.\,.\,.} & $0{,}24357$\\[.4em]
\makebox[17em][l]{$\log\sin(M-u)$\hfill .\,.\,.\,.\,.} & $9{,}40484$\\[.4em]
\makebox[17em][l]{$\log\cos(N-b-P)$\hfill .\,.} & $9{,}86301_n$\\[.4em]
\makebox[17em][l]{$\mathrm{C.}\log\sin P$\hfill .\,.\,.\,.\,.\,.\,.\,.} & $0{,}12099$\\[.4em]
\hline\noalign{\vskip .3em}
\makebox[17em][l]{$\log\left(\frac{\mathrm{d}b}{\mathrm{d}u}\right)$\hfill .\,.\,.\,.\,.\,.\,.\,.\,.\,.\,.} & $9{,}63241_n$\\[.4em]
\end{tabular}
\end{center}
\begin{center}
\begin{tabular}{@{}ll@{}}
\multicolumn{2}{c}{VII$^{*}$}\\[.4em]
\makebox[17em][l]{$\log r\sin u\cos i$\hfill .\,.} & $9{,}75999_n$\\[.4em]
\makebox[17em][l]{$\log\cos(N-b)$\hfill .\,.} & $9{,}99759_n$\\[.4em]
\makebox[17em][l]{$\mathrm{C.}\log\Delta$\hfill .\,.\,.\,.\,.\,.\,.} & $9{,}91759$\\[.4em]
\makebox[17em][l]{$\mathrm{C.}\log\cos N$\hfill .\,.\,.\,.} & $0{,}00001_n$\\[.4em]
\hline\noalign{\vskip .3em}
\makebox[17em][l]{$\log\left(\frac{\mathrm{d}b}{\mathrm{d}i}\right)$\hfill .\,.\,.\,.\,.\,.\,.} & $9{,}67518_n$\\[.4em]
\end{tabular}
\end{center}
\begin{center}
\begin{tabular}{@{}ll@{}}
\multicolumn{2}{c}{VIII}\\[.4em]
\makebox[17em][l]{$({}^{*})$\hfill .\,.\,.\,.\,.\,.\,.\,.\,.\,.\,.} & $9{,}63962$\\[.4em]
\makebox[17em][l]{$\log\sin b\cos b$\hfill .\,.} & $9{,}04212_n$\\[.4em]
\hline\noalign{\vskip .3em}
\makebox[17em][l]{$\log\left(\frac{\mathrm{d}b}{\mathrm{d}\AscendingNode}\right)$\hfill .\,.\,.\,.\,.\,.} & $8{,}68174_n$\\[.4em]
\end{tabular}
\end{center}
\par
Combining these values gives
\[
\begin{aligned}
\mathrm{d}l={}&+0{,}20589\,\mathrm{d}r+1{,}66073\,\mathrm{d}u-0{,}11152\,\mathrm{d}i+1{,}70458\,\mathrm{d}\AscendingNode\\
\mathrm{d}b={}&+0{,}03665\,\mathrm{d}r-0{,}42895\,\mathrm{d}u-0{,}47335\,\mathrm{d}i-0{,}04806\,\mathrm{d}\AscendingNode.
\end{aligned}
\]
We will hardly need to repeat here what we have already pointed out several times. Either the changes $\mathrm{d}l$, $\mathrm{d}b$, $\mathrm{d}u$, $\mathrm{d}i$ and $\mathrm{d}\AscendingNode$ must be expressed in parts of the radius, or the coefficients of $\mathrm{d}r$ itself must be multiplied by $206265''$ if those changes are taken to be expressed in arcseconds.
\par
Denote the longitude of perihelion (in our example $52^0\,18'\,9''{,}30$) by $\Pi$ and the true anomaly by $v$. The longitude in the orbit is then $=u+\AscendingNode=v+\Pi$, so that $\mathrm{d}u=\mathrm{d}v+\mathrm{d}\Pi-\mathrm{d}\AscendingNode$. Substituting this value in the preceding formulas expresses $\mathrm{d}l$ and $\mathrm{d}b$ in terms of $\mathrm{d}r$, $\mathrm{d}v$, $\mathrm{d}\Pi$, $\mathrm{d}\AscendingNode$ and $\mathrm{d}i$. \SourceLeaf{104}It therefore remains only to express $\mathrm{d}r$ and $\mathrm{d}v$ in terms of the differential changes in the elliptic elements, following Arts. 15 and 16\footnote{Everyone will readily see that the symbol $M$ in the following calculation ceases to denote our auxiliary angle and again denotes the mean anomaly, as in Section I.}.
\par
In our example, in Art. 14, we had $\log\frac{r}{a}=9{,}90355=\log\left(\frac{\mathrm{d}r}{\mathrm{d}a}\right)$.
\begin{center}
\begin{tabular}{@{}ll@{}}
\makebox[17em][l]{$\log\frac{aa}{rr}$\hfill .\,.\,.\,.\,.\,.} & $0{,}19290$\\[.4em]
\makebox[17em][l]{$\log\cos\varphi$\hfill .\,.\,.\,.} & $9{,}98652$\\[.4em]
\hline\noalign{\vskip .3em}
\makebox[17em][l]{$\log\left(\frac{\mathrm{d}v}{\mathrm{d}M}\right)$\hfill .\,.\,.\,.} & $0{,}17942$\\[.4em]
\end{tabular}
\end{center}
\begin{center}
\begin{tabular}{@{}ll@{}}
$2-e\cos E=$ & $1{,}80085$\\[.4em]
$ee=$ & $0{,}06018$\\[.4em]
\hline\noalign{\vskip .3em}
 & $1{,}74067$\\[.4em]
\makebox[17em][l]{$\log$\hfill .\,.\,.\,.\,.\,.\,.\,.} & $0{,}24072$\\[.4em]
\makebox[17em][l]{$\log\frac{aa}{rr}$\hfill .\,.\,.\,.\,.\,.} & $0{,}19290$\\[.4em]
\makebox[17em][l]{$\log\sin E$\hfill .\,.\,.\,.} & $9{,}76634_n$\\[.4em]
\hline\noalign{\vskip .3em}
\makebox[17em][l]{$\log\left(\frac{\mathrm{d}v}{\mathrm{d}\varphi}\right)$\hfill .\,.\,.\,.\,.} & $0{,}19996_n$\\[.4em]
\end{tabular}
\end{center}
\begin{center}
\begin{tabular}{@{}ll@{}}
\makebox[17em][l]{$\log a$\hfill .\,.\,.\,.\,.\,.\,.} & $0{,}42244$\\[.4em]
\makebox[17em][l]{$\log\operatorname{tang}\varphi$\hfill .\,.\,.} & $9{,}40320$\\[.4em]
\makebox[17em][l]{$\log\sin v$\hfill .\,.\,.\,.\,.} & $9{,}84931_n$\\[.4em]
\hline\noalign{\vskip .3em}
\makebox[17em][l]{$\log\left(\frac{\mathrm{d}r}{\mathrm{d}M}\right)$\hfill .\,.\,.\,.} & $9{,}67495_n$\\[.4em]
\end{tabular}
\end{center}
\begin{center}
\begin{tabular}{@{}ll@{}}
\makebox[17em][l]{$\log a$\hfill .\,.\,.\,.\,.\,.\,.} & $0{,}42244$\\[.4em]
\makebox[17em][l]{$\log\cos\varphi$\hfill .\,.\,.\,.} & $9{,}98652$\\[.4em]
\makebox[17em][l]{$\log\cos v$\hfill .\,.\,.\,.} & $9{,}84966$\\[.4em]
\hline\noalign{\vskip .3em}
\makebox[17em][l]{$\log\left(\frac{\mathrm{d}r}{\mathrm{d}\varphi}\right)$\hfill .\,.\,.\,.\,.} & $0{,}25862_n$\\[.4em]
\end{tabular}
\end{center}
\noindent From this we obtain
\[
\begin{aligned}
\mathrm{d}v&=+1{,}51154\,\mathrm{d}M-1{,}58475\,\mathrm{d}\varphi\\
\mathrm{d}r&=-0{,}47310\,\mathrm{d}M-1{,}81393\,\mathrm{d}\varphi+0{,}80085\,\mathrm{d}a,
\end{aligned}
\]
and substituting these values in the preceding formulas gives
\[
\begin{aligned}
\mathrm{d}l={}&+2{,}41287\,\mathrm{d}M-3{,}00531\,\mathrm{d}\varphi+0{,}16488\,\mathrm{d}a\\
 &+1{,}66073\,\mathrm{d}\Pi-0{,}11152\,\mathrm{d}i+0{,}04385\,\mathrm{d}\AscendingNode\\
\mathrm{d}b={}&-0{,}66572\,\mathrm{d}M+0{,}61331\,\mathrm{d}\varphi+0{,}02935\,\mathrm{d}a\\
 &-0{,}42895\,\mathrm{d}\Pi-0{,}47335\,\mathrm{d}i+0{,}38089\,\mathrm{d}\AscendingNode.
\end{aligned}
\]
\par
Suppose that the time corresponding to the computed position is $n$ days from the epoch, and denote the mean longitude at the epoch by $N$ and the daily motion by $\DailyMotionSign$. Then $M=N+n\DailyMotionSign-\Pi$, and hence $\mathrm{d}M=\mathrm{d}N+n\,\mathrm{d}\DailyMotionSign-\mathrm{d}\Pi$. In our example, the time corresponding to the computed position is October $17{,}41501$, 1804, reckoned at the Paris meridian. If, then, the beginning of 1805 is taken as the epoch, $n=-74{,}58499$. The mean longitude established for that epoch \SourceLeaf{105}was $=41^0\,52'\,21''{,}68$ and the daily motion $=824''{,}7990$. Substituting the value of $\mathrm{d}M$ in the formulas just obtained gives the following differential changes in the geocentric position in terms of changes in the elements alone.
\[
\begin{aligned}
\mathrm{d}l={}&+2{,}41287\,\mathrm{d}N-179{,}96\,\mathrm{d}\DailyMotionSign-0{,}75214\,\mathrm{d}\Pi\\
 &-3{,}00531\,\mathrm{d}\varphi+0{,}16488\,\mathrm{d}a-0{,}11152\,\mathrm{d}i+0{,}04385\,\mathrm{d}\AscendingNode\\
\mathrm{d}b={}&-0{,}66572\,\mathrm{d}N+49{,}65\,\mathrm{d}\DailyMotionSign+0{,}23677\,\mathrm{d}\Pi\\
 &+0{,}61331\,\mathrm{d}\varphi+0{,}02935\,\mathrm{d}a-0{,}47335\,\mathrm{d}i+0{,}38089\,\mathrm{d}\AscendingNode.
\end{aligned}
\]
\par
If the mass of the celestial body is either neglected or at least regarded as known, $\DailyMotionSign$ and $a$ will depend on each other, so either $\mathrm{d}\DailyMotionSign$ or $\mathrm{d}a$ can be eliminated from our formulas. Indeed, Art. 6 gives $\DailyMotionSign a^{\frac{3}{2}}=k\BareRoot{(1+\mu)}$, so we have $\frac{\mathrm{d}\DailyMotionSign}{\DailyMotionSign}=-\frac{3}{2}\frac{\mathrm{d}a}{a}$. In this formula, if $\mathrm{d}\DailyMotionSign$ is to be expressed in parts of the radius, $\DailyMotionSign$ must be expressed in the same way. Thus in our example we have
\begin{center}
\begin{tabular}{@{}ll@{}}
\makebox[17em][l]{$\log\DailyMotionSign$\hfill .\,.\,.\,.\,.} & $2{,}91635$\\[.4em]
\makebox[17em][l]{$\log1''$\hfill .\,.\,.\,.} & $4{,}68557$\\[.4em]
\makebox[17em][l]{$\log\frac{3}{2}$\hfill .\,.\,.\,.\,.} & $0{,}17609$\\[.4em]
\makebox[17em][l]{$\mathrm{C.}\log a$\hfill .\,.\,.} & $9{,}57756$\\[.4em]
\hline\noalign{\vskip .3em}
\makebox[17em][l]{$\log\frac{\mathrm{d}\DailyMotionSign}{\mathrm{d}a}$\hfill .\,.\,.\,.} & $7{,}35557_n$\\[.4em]
\end{tabular}
\end{center}
\noindent or $\mathrm{d}\DailyMotionSign=-0{,}0022676\,\mathrm{d}a$, and $\mathrm{d}a=-440{,}99\,\mathrm{d}\DailyMotionSign$. Substituting this value in our formulas brings us to their final form.
\[
\begin{aligned}
\mathrm{d}l={}&+2{,}41287\,\mathrm{d}N-252{,}67\,\mathrm{d}\DailyMotionSign-0{,}75214\,\mathrm{d}\Pi\\
 &-3{,}00531\,\mathrm{d}\varphi-0{,}11152\,\mathrm{d}i+0{,}04385\,\mathrm{d}\AscendingNode\\
\mathrm{d}b={}&-0{,}66572\,\mathrm{d}N+36{,}71\,\mathrm{d}\DailyMotionSign+0{,}23677\,\mathrm{d}\Pi\\
 &+0{,}61331\,\mathrm{d}\varphi-0{,}47335\,\mathrm{d}i+0{,}38089\,\mathrm{d}\AscendingNode.
\end{aligned}
\]
In deriving these formulas we assumed that all the changes $\mathrm{d}l$, $\mathrm{d}b$, $\mathrm{d}N$, $\mathrm{d}\DailyMotionSign$, $\mathrm{d}\Pi$, $\mathrm{d}\varphi$, $\mathrm{d}i$ and $\mathrm{d}\AscendingNode$ were expressed in parts of the radius. By the homogeneity of all the terms, however, the same formulas clearly still hold if all those changes are expressed in arcseconds.
\begin{center}\rule{7em}{.4pt}\end{center}
